\documentclass[UTF8,fontset=none,12pt]{ctexart}

\setCJKmainfont{Songti SC}
\setCJKsansfont{PingFang SC}
\setCJKmonofont{Songti SC}

\usepackage[a4paper,margin=2.35cm]{geometry}
\usepackage{amsmath,amssymb,amsthm,mathtools,bm}
\usepackage{graphicx}
\usepackage{float}
\usepackage{booktabs}
\usepackage{enumitem}
\usepackage[hidelinks]{hyperref}

\newcommand{\C}{\mathbb{C}}
\newcommand{\R}{\mathbb{R}}
\newcommand{\F}{\mathbb{F}}
\newcommand{\fl}{\operatorname{fl}}
\newcommand{\epsmach}{\varepsilon_{\mathrm{machine}}}
\theoremstyle{definition}
\newtheorem{theorem}{定理}
\renewcommand{\thetheorem}{13.\arabic{theorem}}
\renewcommand{\thefigure}{13.\arabic{figure}}

\title{数值线性代数\quad 第 13 讲：浮点运算}
\author{Lloyd N. Trefethen、David Bau III\quad 中文阅读版}
\date{}

\begin{document}
\maketitle
\pagestyle{plain}

计算机发明后不久，人们就对数字设备表示实数的正确方式形成了共识.秘密就在于浮点运算：它是科学记数法在硬件中的对应物.在开始研究数值线性代数算法的准确性之前，必须先考察这一主题.

\section*{数字表示的局限}

由于数字计算机使用有限个比特表示实数，它们只能表示实数（或者复数；本讲末尾会讨论复数）的一个有限子集.这一局限带来两个困难.第一，被表示的数不可能任意大或任意小.第二，相邻表示数之间必然存在间隔.

现代计算机能够表示足够大和足够小的数，因此第一个限制很少造成困难.例如，广泛使用的 IEEE 双精度运算允许表示最大约为 $1.79\times10^{308}$、最小约为 $2.23\times10^{-308}$ 的数；对本书讨论的大多数问题而言，这个范围已经足够大.换言之，溢出和下溢通常不是严重风险（但如果有人要求你计算行列式，就要特别小心）.

相反，表示数之间的间隔在科学计算中始终值得关注.例如，在 IEEE 双精度运算中，区间 $[1,2]$ 表示为离散子集
\begin{equation}
    1,\quad 1+2^{-52},\quad 1+2\cdot2^{-52},\quad
    1+3\cdot2^{-52},\quad\ldots,\quad 2.
    \tag{13.1}
\end{equation}

区间 $[2,4]$ 则由同一组数乘以 2 得到：
\[
    2,\quad 2+2^{-51},\quad 2+2\cdot2^{-51},\quad
    2+3\cdot2^{-51},\quad\ldots,\quad 4.
\]
一般地，区间 $[2^j,2^{j+1}]$ 由式 (13.1) 中的数乘以 $2^j$ 得到.因此，在 IEEE 双精度运算中，相邻表示数之间的间隔从相对意义上看从不会大于
\[
    2^{-52}\approx2.22\times10^{-16}.
\]
如果使用稳定算法，这个量对大多数目的而言确实小得可以忽略（见下一讲）.但令人意外的是，有多少粗心设计的算法最终会表现出不稳定性!

\section*{浮点数}

IEEE 运算是以浮点表示实数的一种运算系统.如今，通用计算机普遍采用这种做法.在浮点数系统中，小数点（或二进制点）的位置与数字本身分开存储；相邻表示数之间的间隔会随数的大小成比例变化.这与定点表示不同，后者的间隔始终相同.

具体地，考虑如下理想化的浮点数系统.令 $\F$ 是实数 $\R$ 的一个离散子集，它由整数 $\beta\ge2$ 和整数 $t\ge1$ 决定；$\beta$ 称为基数或底（通常取 2），$t$ 称为精度（IEEE 单精度和双精度分别对应 24 和 53）. $\F$ 中包含数字 0，以及所有形如
\begin{equation}
    x=\pm\left(\frac{m}{\beta^t}\right)\beta^e
    \tag{13.2}
\end{equation}
的数，其中 $m$ 是满足 $1\le m\le\beta^t$ 的整数，$e$ 是任意整数.

等价地，也可以限制
\[
    \beta^{t-1}\le m\le\beta^t-1,
\]
这样 $m$ 的选择就是唯一的.此时 $\pm(m/\beta^t)$ 称为 $x$ 的小数部分或尾数，$e$ 称为指数.

上述浮点数系统是理想化的，因为它忽略了上溢和下溢.因此，$\F$ 是一个可数无穷集，并且具有自相似性：
\[
    \F=\beta\F.
\]

\section*{机器 epsilon}

$\F$ 的分辨率通常用一个称为机器 epsilon（machine epsilon）的数来概括.暂时定义
\begin{equation}
    \epsmach=\frac12\beta^{1-t}.
    \tag{13.3}
\end{equation}
（在式 (13.7) 之后我们会修改这个定义.）这个数是 1 与下一个更大的浮点数之间距离的一半.从相对意义上看，浮点数之间的间隔最大也就是这个量.也就是说，$\epsmach$ 具有如下性质：
\begin{equation}
    \text{对所有 }x\in\R,\ \text{存在 }x'\in\F,\quad
    |x-x'|\le\epsmach|x|.
    \tag{13.4}
\end{equation}

对于各种计算机上常见的 $\beta$ 和 $t$，$\epsmach$ 通常位于 $10^{-6}$ 与 $10^{-35}$ 之间.在 IEEE 单精度和双精度运算中，分别规定
\[
    \epsmach=2^{-24}\approx5.96\times10^{-8},
    \qquad
    \epsmach=2^{-53}\approx1.11\times10^{-16}.
\]

令 $\fl:\R\to\F$ 表示把实数取为浮点系统中最近浮点数的函数，即舍入后的等价值（平局情形可以任意处理；如何处理平局以避免统计偏差，本身也是一个有趣的问题）.式 (13.4) 可以用 $\fl$ 表述为：
\begin{equation}
    \text{对所有 }x\in\R,\ \text{存在 }\delta,\ |\delta|\le\epsmach,
    \quad \fl(x)=x(1+\delta).
    \tag{13.5}
\end{equation}
也就是说，实数与其最近浮点近似之间的差，从相对意义上看总小于 $\epsmach$.

\section*{浮点运算}

当然，仅仅表示实数是不够的；还必须对它们进行计算.在计算机上，所有数学计算都归结为若干基本算术运算，经典的运算集合是 $+$、$-$、$\times$ 和 $\div$.从数学上说，这些符号表示 $\R$ 上的运算.在计算机上，它们有作用于 $\F$ 的对应运算.通常把这些浮点运算记为 $\oplus$、$\ominus$、$\otimes$ 和 $\oslash$.

计算机可以按照下面的设计原则构造.令 $x,y\in\F$ 为任意浮点数.令 $*$ 为 $+$、$-$、$\times$ 或 $\div$ 中的一个运算，并令 $\circledast$ 为它对应的浮点运算.那么要求
\begin{equation}
    x\circledast y=\fl(x*y).
    \tag{13.6}
\end{equation}
如果这一性质成立，那么由式 (13.5) 和 (13.6) 可以得到一个简单而有力的性质.

\begin{center}
\fbox{\begin{minipage}{0.92\linewidth}
\textbf{浮点运算基本公理}

对所有 $x,y\in\F$，存在 $\delta$ 满足 $|\delta|\le\epsmach$，使得
\[
    x\circledast y=(x*y)(1+\delta).
\tag{13.7}
\]
\end{minipage}}
\end{center}

换言之，每一次浮点运算都精确到一个至多为 $\epsmach$ 的相对误差.

\section*{再次讨论机器 epsilon}

本书的舍入误差分析建立在式 (13.5) 和 (13.7) 上，而不是建立在前面描述的浮点运算其他细节上.这意味着，我们可以宽容地考虑一些并不完全按照式 (13.6) 完美执行浮点计算的硬件实现.对于这样的机器，只要把 $\epsmach$ 换成一个稍大的值，式 (13.5) 和 (13.7) 仍可能成立.例如，在中间量采用截断而不是舍入的计算机上，式 (13.7) 可能需要把 $\epsmach$ 换成 $2\epsmach$.

处理这类复杂情况的最简单方式，是保留式 (13.5) 和 (13.7) 的形式，但修改 $\epsmach$ 的定义.从现在开始，假设 $\epsmach$ 定义为使式 (13.5) 和 (13.7) 成立的最小数.对于包括所有采用 IEEE 运算的计算机在内的大多数计算机，这一修改不会显著改变其数值.

有时，为使式 (13.7) 成立，可能需要一个出乎意料地大的 $\epsmach$.1994 年底，Intel Pentium\textsuperscript{\texttrademark} 微处理器曾因 IEEE 双精度标准实现中的查表错误而声名狼藉；其有效精度粗糙了 11 个数量级，$\epsmach\approx6.1\times10^{-5}$（这个错误很快得到纠正）.事实上，有些机器上式 (13.7) 只有在 $\epsmach=1$ 时才成立.例如，直到 20 世纪 90 年代中期，Cray 计算机上的浮点减法都具有这一性质，因为减法实现时没有“保护数字”（guard digit）.这些计算机并非毫无用处，但它们要求采用不同于本书的误差分析方式.

幸运的是，式 (13.7) 公理的优点，以及采用统一计算机算术标准的好处，近年来已经被计算机制造商广泛接受.市场上不满足“小 $\epsmach$”版本式 (13.7) 的机器正在减少.事实上，IEEE 运算本身正在迅速成为各种规模计算机的标准；截至 1996 年，它已经用于所有 IBM 兼容个人计算机，以及 SUN、DEC、Hewlett--Packard 和 IBM 制造的所有工作站.

\section*{复数浮点运算}

浮点复数通常表示为一对浮点实数，作用于它们的基本运算则归结为对实部和虚部进行运算.结果是，式 (13.7) 对复浮点数同样成立，但乘法和除法例外：对于 $\otimes$ 和 $\oslash$，必须把式 (13.3) 中的 $\epsmach$ 分别放大到原来量级的
\[
    2^{3/2}\quad\text{和}\quad 2^{5/2}.
\]
一旦按这种方式调整 $\epsmach$，复数的舍入误差分析就可以像实数一样进行.

\section*{习题}

\begin{enumerate}[label=\textbf{13.\arabic*}]
    \item 在一对相邻的非零 IEEE 单精度实数之间，有多少个 IEEE 双精度数？

    \item 式 (13.2) 定义的浮点数系统 $\F$ 包含许多整数，但并不包含全部整数.
    \begin{enumerate}[label=(\alph*)]
        \item 给出不属于 $\F$ 的最小正整数 $n$ 的精确公式.
        \item 对 IEEE 单精度和双精度运算，$n$ 的值分别是多少？
        \item 设法在自己的计算机上验证这一结果.具体地，设计并运行一个程序，给出证据说明 $n-3$、$n-2$ 和 $n-1$ 属于 $\F$，而 $n$ 不属于 $\F$. $n+1$、$n+2$ 和 $n+3$ 呢？
    \end{enumerate}

    \item 考虑多项式
    \begin{equation*}
    \begin{aligned}
        p(x)&=(x-2)^9\\
        &=x^9-18x^8+144x^7-672x^6+2016x^5\\
        &\quad-4032x^4+5376x^3-4608x^2+2304x-512.
    \end{aligned}
    \end{equation*}
    \begin{enumerate}[label=(\alph*)]
        \item 对 $x=-1.920,-1.919,-1.918,\ldots,2.080$ 作图；通过系数 $1,-18,144,\ldots$ 计算 $p(x)$.
        \item 再作一次同样的图，但这次使用表达式 $(x-2)^9$ 计算 $p(x)$.
    \end{enumerate}

    \item 多项式
    \[
        p(x)=x^5-2x^4-3x^3+3x^2-2x-1
    \]
    有三个实零点.对 $p$ 使用初值 $x_0=0$ 的 Newton 方法，会产生快速收敛到某个零点 $x_*\approx-0.315$ 的估计序列 $x_1,x_2,x_3,\ldots$.
    \begin{enumerate}[label=(\alph*)]
        \item 在 $\epsmach\approx10^{-16}$ 的浮点运算中计算 $x_1,\ldots,x_6$.估计这些数中每一个有多少位数字正确？
        \item 借助 MAPLE 或 MATHEMATICA 等符号代数系统，再次精确计算 $x_1,\ldots,x_6$.每个 $x_j$ 都是有理数.对于每个 $j$，分子和分母各有多少位数字？
    \end{enumerate}
\end{enumerate}

\end{document}
